\begin{textsample}{POS Dim 8 – global\_north\_2025\_02 – Score 16.00 – global\_north\_2025\_02/120916419.txt}  \label{ex:f8_pos_003}
Follow us Palestinians defy Trump ' s \textbf{ethnic} \textbf{cleansing} plan and return home Despite relentless attacks and international pressure, Palestinians in Gaza are courageously reclaiming their homes, rejecting attempts to erase their existence A young woman wears a dress in the colours of the Palestinian flag as people return to Rafah in the southern Gaza Strip on 19 January, 2025 ( AFP ) Since assuming office, US President Donald Trump has relentlessly urged Egypt, Jordan, and other Muslim-majority countries to \textbf{resettle} Palestinians from Gaza.
Although Palestinians have firmly rejected Trump ' s proposal, it has continued to dominate the front pages of almost every Israeli newspaper.
Israel ' s Finance Minister Bezalel Smotrich, who last year argued that it was"justified and moral"to starve Palestinians in Gaza, has been outspoken in his support of the idea, stating :"After 76 years in which most of Gaza ' s population was forcibly held in harsh conditions to preserve the aspiration to destroy the State of Israel, the idea of helping them @ @ @ @ @ @ @ @ @ @ is a great one."Yedioth Ahronoth ' s senior military correspondent, Yossi Yehoshua, has also been a staunch supporter, suggesting :"Perhaps the time has come to adopt Trump ' s proposal and discuss voluntary exile from Gaza."On Tuesday, at a joint press conference alongside Israeli Prime Minister Benjamin Netanyahu, Trump went a step further and announced the US will be taking over and running Gaza, potentially for the foreseeable future.
New \textbf{MEE} \textbf{newsletter} : Jerusalem Dispatch Sign up to get the latest \textbf{insights} and analysis on Israel-Palestine, alongside Turkey \textbf{Unpacked} and other \textbf{MEE} \textbf{newsletters} Shortly after Israel launched its war on Gaza in October 2023, fears quickly emerged that Israel would execute its undeclared plan to ethnically cleanse Palestinians from the enclave.
Given the high level of support Israel was receiving from its western backers, many of us feared that a similar fate would await Palestinians in Jerusalem, the West Bank, and, eventually, even those of us living in the lands of historic Palestine @ @ @ @ @ @ @ @ @ @ a genuine belief that, no matter how dire the situation becomes, the Palestinian people will not disappear This concern stemmed from a 10-page document issued in October 2023 by Israeli Minister Gila Gamliel ' s Intelligence Ministry, which proposed forcibly transferring Palestinians from Gaza to Egypt ' s Sinai Peninsula.
Gamliel ' s document outlined three alternatives for post-war Gaza, with the option"that will yield positive, long-term strategic results"involving the expulsion of Palestinians to Sinai."We affirm our rejection of any attempts to compromise Palestinians ' unalienable rights, whether through settlement activities, or evictions or annexation of land or through vacating the land from its owners,"they said in a joint statement.
Even Trump ' s"favourite dictator", Egyptian President Abdel Fattah el-Sisi, has voiced dissent, warning that Egyptians would take to the streets to express their disapproval.
Palestinian defiance Amid all the talk of \textbf{ethnic} \textbf{cleansing}, Palestinians have remained resolute, with \textbf{extraordinary} scenes unfolding in northern Gaza.
The image of @ @ @ @ @ @ @ @ @ @ Gaza after being displaced in the south evokes memories of the Nakba, when hundreds of thousands were forced to flee their homes due to Zionist militias and armed gangs.
But this time, the scene and mood are not ones of despair.
There is now a genuine belief that, no matter how dire the situation becomes, the Palestinian people will not disappear. ' Clean out ' Gaza : Why Trump ' s ' voluntary ' transfer proposal should be made to Israelis As a result, Israeli media has gone into a complete meltdown, with many lamenting the scenes of Palestinian defiance.
Channel 13 ' s political correspondent, Moriah Asraf, recently expressed :"These images make me shiver all over my body...
Something about the Gazans returning to their homes, albeit destroyed, but to their homes - it drives me crazy."Matan Zuri, a security correspondent for Ynet, wrote :"Thousands of Palestinians have returned to the devastated northern Gaza Strip.
The dream of renewed Jewish settlement @ @ @ @ @ @ @ @ @ @ price of ending the war and returning the hostages.
We knew it would happen, we saw it coming, and there was no choice but to accept it with submission and stick to the goodness of the deal."Whatever happens next is anyone ' s guess, but the image of Palestinians returning to what remains of their bombed-out homes has been the most powerful response yet to Trump ' s racist and dehumanising plan.
The views expressed in this article belong to the author and do not \textbf{necessarily} reflect the editorial policy of Middle East Eye.
Lubna Masarwa is a journalist and Middle East Eye ' s Palestine and Israel bureau chief, based in Jerusalem.
Middle East Eye delivers independent and \textbf{unrivalled} \textbf{coverage} and analysis of the Middle East, North Africa and beyond.
To learn more about \textbf{republishing} this content and the \textbf{associated} \textbf{fees}, \textbf{please} \textbf{fill} out this form.
More about \textbf{MEE} can be found here.

% matched lemmas: associate, cleansing, coverage, ethnic, extraordinary, fee, fill, insight, mee, necessarily, newsletter, please, republish, resettle, unpack, unrivalled
\end{textsample}
